\clearpage\fancyhead[L]{\small JUNO.}\hypertarget{native-377-lower}{}\label{native-377-lower}
\begin{center}\textbf{JUNO.}\par\vspace{4pt}\hrule\vspace{4pt}\fontsize{10}{12}\selectfont Göttingische gelehrte Anzeigen. No. 204. pp. 2033\;..\;2034. 23 December 1813.\par\vspace{4pt}\hrule\vspace{4pt}\end{center}


\par\indent The season in which \emph{Juno} came into opposition with the Sun in \emph{George’s Harp} this year is usually unfavorable for astronomical observations here. Professor \textsc{Gauss} was able to obtain only one observation of this planet around that time:
\begin{center}\fontsize{10}{12}\selectfont\setlength{\tabcolsep}{4pt}\renewcommand{\arraystretch}{1.15}\begin{tabular}{l|c|c|c}
1813 & Mean time & Right ascension & Declination\\\hline
Nov. 19 & $9^{\mathrm u}\;46^{\mathrm m}\;46''$ & $60^\circ\;30'\;35''0$ & $3^\circ\;5'\;9''3\ \text{S.}$\\
\end{tabular}\end{center}
\par\indent Although the planet was also very close to its perihelion this time, and so was in the most favorable position for illumination, it barely reached the 8th magnitude.

\Needspace{14\baselineskip}\par\indent There are no observations from elsewhere, and even if the planet was observed somewhere else, none can be expected at present from distant observatories. Professor \textsc{Gauss} therefore had to decide to base the calculation of the opposition on the individual observations above. The result is as follows:
\begin{center}\leavevmode\SourceLeaf{378}\emph{Eighth opposition of Juno.}\end{center}
\begin{center}1813 November 19. $18^{\mathrm u}\;11^{\mathrm m}\;5^{\mathrm s}$ Mean time in Göttingen.\end{center}
\begin{center}\fontsize{10}{12}\selectfont\setlength{\tabcolsep}{2pt}\renewcommand{\arraystretch}{1.15}\begin{tabular}{lcr}
True longitude & $.\;.\;.\;.\;.\;.\;.\;.\;.\;.\;.\;.\;.$ & $57^\circ\;33'\;58''0$\\
Geocentric latitude & $.\;.\;.\;.\;.\;.\;.\;.\;.\;.\;.$ & $23\;18\;46.8\ \text{S.}$\\
\end{tabular}\end{center}
\Needspace{20\baselineskip}\par\indent For an ephemeris to be calculated for the planet’s next period of visibility and published elsewhere, Professor \textsc{Gauss} derived new elliptic elements from the fifth, sixth, seventh and eighth oppositions.
\begin{center}\fontsize{9}{11}\selectfont\setlength{\tabcolsep}{2pt}\renewcommand{\arraystretch}{1.15}\begin{tabular}{lcr}
\shortstack[l]{Epoch of mean longitude 1810\\at the Göttingen meridian} &  & $95^\circ\;29'\;56''$\\
Daily tropical mean motion & $.\;.\;.\;.\;.\;.\;.$ & $812''709$\\
Tropical orbital period & $.\;.\;.\;.\;.\;.\;.\;.\;.\;.\;.\;.\;.$ & $1594\ \text{days}\ 16\ \text{hours}$\\
Longitude of perihelion (1810) & $.\;.\;.\;.\;.\;.\;.\;.\;.\;.\;.$ & $53^\circ\;6'\;40''$\\
Longitude of ascending node (1810) & $.\;.\;.\;.\;.\;.\;.$ & $171\;7\;13$\\
Angle of eccentricity & $.\;.\;.\;.\;.\;.\;.\;.\;.\;.\;.\;.\;.\;.$ & $14\;43\;8$\\
Inclination of the orbit & $.\;.\;.\;.\;.\;.\;.\;.\;.\;.\;.\;.\;.\;.$ & $13\;4\;12$\\
Log. of the semimajor axis & $.\;.\;.\;.\;.\;.\;.\;.$ & $0.4267670$\\
\end{tabular}\end{center}
\par\indent Comparison of these elements with those based on the \emph{first} four oppositions (see no. 136 of this journal for 1808 [p. 306 of this volume]) shows very clearly how strongly perturbations from other bodies affect this planet too. The difference in the mean motion is particularly striking. The earlier oppositions gave an orbital period shorter by 3 days and 4 hours. Without the perturbations this difference could not have amounted to half an hour.
